\subsection{Application of Formula Selecting Networks}

Formula selecting networks are efficient in the simultaneous representation of multiple formulas with similar structure.


%There are two main applications of formula selecting networks.
%First, when contracting the selection variables with a weight tensor we get a weighted sum of the parametrized formulas.
%Second, when contracting the categorical variables with a distribution or a knowledge base, we get a tensor storing the satisfaction rates respectively the world counts of the parametrized formulas.

%\subsubsection{Representation of selection encodings}
%
%The main application of formula selecting networks in this work is the efficient representation of selection encodings.
%This will be exploited in the sparse representation of exponential families by energies and in structure learning.
%In the next lemma we will show the correspondence of formula selecting networks and selection encodings.
%
%\begin{lemma}
%    \label{lem:relToSelFSN}
%    Given a set $\{\formulaof{\selindexlist} : \selindexlist\in\selstates\}$ of propositional formulas we define the statistic
%    \[ \formulaset : \catindices \rightarrow (\formulaof{\selindexlist}(\catindices))_{\selindexlist} \, . \]
%    and the formula selecting map
%    \[ \fselectionmap: \catindices , \selindexlist \rightarrow \formulaof{\selindexlist} (\catindices) \, . \]
%    Then
%    \[ \sencodingofat{\formulaset}{\shortcatvariablelist, \shortselvariablelist} = \fselectionmap\left[\shortcatvariablelist, \shortselvariablelist \right] \, .  \]
%\end{lemma}
%\begin{proof}
%    For any indices $\shortselindices\in\selstates$ and $\shortcatindices\in\atomstates$ we have
%    \begin{align*}
%        \sencodingofat{\formulaset}{\shortcatvariablelist=\shortcatindices, \shortselvariablelist=\shortselindices}
%        =  \formulaof{\selindexlist}(\catindices) =  \fselectionmap\left[\shortcatvariablelist=\shortcatindices, \shortselvariablelist=\shortselindices \right] \, .
%    \end{align*}
%\end{proof}
%
%%% Reason for basis encodings and selection encodings.
%Technically, basis encodings have been exploited to derive decompositions based on basis calculus.
%Selection encodings on the other hand enable the application of formula selecting networks as superpositions of formulas.



\subsubsection{Efficient Representation of Formulas}

% Exponentially many formulas represented by linear demand
Formula Selecting Neural Networks are means to represent exponentially many formulas with linear (in sum of candidates list lengths) storage.
Their contraction with probability tensor networks, is thus a batchwise evaluation of exponentially many formulas.
This is possible due to redundancies in logical calculus due to modular combinations of subformulas.

% Retrieve functions
We can retrieve specific formulas by slicing the selection variables, i.e. for $\selindices$ we have
\[ \exformula_{\selindices}[\shortcatvariables] = \sencfselectionmapat{\shortcatvariables,\selvariable=\selindices} \, .  \]

In a tensor network diagram we depict this by
\begin{center}
    \begin{tikzpicture}[thick, scale=0.35] % , baseline = -3.5pt

    \node[anchor=east] (text) at (-3,0) {${\exformula_{\selindices}} \quad\quad {=}$};

    \drawatomindices{0}{-4}
    \draw (-1,3) rectangle (5, -3);
    \node[anchor=center] (text) at (2,0) {$\bsencodingof{\fselectionmap}$};

    \draw[->-] (2,3)--(2,4.5);
    \drawvariabledot{2}{4}
    \draw[] (2,5)--(2,4) node[right] {\colorlabelsize $\headvariableof{\fselectionmap}$};
    \draw (1,5) rectangle (3,7);
    \node[anchor=center] (text) at (2,6) {$\tbasis$};

    \draw[] (7,-2.25)--(6,-2.25) node[below] {\colorlabelsize $\selvariableof{0}$};
    \drawvariabledot{6}{-2.25}
    \draw[-<-] (5,-2.25)--(6.5,-2.25);

    \draw[] (7,-0.5)--(6,-0.5) node[below] {\colorlabelsize $\selvariableof{1}$};
    \drawvariabledot{6}{-0.5}
    \draw[-<-] (5,-0.5)--(6.5,-0.5);


    \node[anchor=center] (text) at (6,1) {$\vdots$};
    \draw[] (7,2.25)--(6,2.25) node[above] {\colorlabelsize $\selvariableof{\selorder\shortminus1}$};
    \drawvariabledot{6}{2.25}
    \draw[-<-] (5,2.25)--(6.5,2.25);

    \draw (7,1.5) rectangle (9,3);
    \node[anchor=center] (text) at (8,2.25) {$\onehotmapof{\selindexof{\selorder\shortminus1}}$};

    \draw (7,0.25) rectangle (9,-1.25);
    \node[anchor=center] (text) at (8,-0.5) {$\onehotmapof{\selindexof{1}}$};

    \draw (7,-1.5) rectangle (9,-3);
    \node[anchor=center] (text) at (8,-2.25) {$\onehotmapof{\selindexof{0}}$};

\end{tikzpicture}
\end{center}

% Interpretation by dynamic programming
Another perspective on the efficient formula evaluation by selection tensor networks is dynamic computing.
Evaluating a formula requires evaluations of its subformulas, which are done by subcontractions and saved for different subformulas due to the additional selection legs.

% Storage problem of solutions
However, we need to avoid contracting the tensor with leaving all selection legs open, since this would require exponential storage demand.
% Sparse algorithm
We can avoid this storage bottleneck by contraction of parameter cores $\canparam$ with efficient network decompositions along the selection variables. %extending the contractions by additional cores leaving less variable legs open.

% Gibbs sampling
In Gibbs Sampling (\algoref{alg:Gibbs}), one can use the energy-based approach to queries \theref{the:energyContractionQueries}, and contract basis vectors on all but one selection variables.

%\red{This is the case when contracting gradients of the parameter tensor networks in alternating least squares approaches.
%Other methods avoiding the bottleneck can be constructed by MCMC sampling, for example Gibbs Sampling.
%Here we only need to vary local components of the formula reflected in keeping only single variable legs open.}



\subsubsection{Batch Contraction of Parametrized Formulas}

Given a set $\formulaset$ of formulas, we build a formula selecting network parametrizing the formulas.
The contraction
\begin{align*}
    \contractionof{\extnet,\fselectionmap}{\shortselvariables}
\end{align*}
is a tensor containing the contractions of the formulas $\formulaof{\shortselindices}$ with an arbitrary tensor network $\extnet$ as
\begin{align*}
    \contraction{\extnet,\formulaof{\shortselindices}} = \contractionof{\extnet,\fselectionmap}{\shortselvariables=\shortselindices} \, .
\end{align*}


\subsubsection{Average Contraction of Parametrized Formulas}

We show in the next two examples, how a full contraction of the formula selecting map with a probability distribution or a knowledge base can be interpreted.

\begin{example}[Average satisfaction of formulas]
    The average of the formula satisfactions in $\formulaset$ given a probability tensor $\probtensor$ is
    \[ \frac{1}{\prod_{\selenumeratorin}\seldimof{\selenumerator}} \cdot \contraction{\probtensor,\sencodingof{\formulaset}} \, . \]
\end{example}


\begin{example}[Deciding whether any formula is not contradicted]
    For example: We want to decide, whether there is a formula in $\formulaset$ not contradicted by a Knowledge base $\kb$.
    This is the case if and only if
    \[ \contraction{\kb,\sencodingof{\formulaset}} = 0 \, .  \]
%    We use \lemref{lem:relToSelFSN} to get that $\sencodingof{\formulaset}=\fselectionmap$.
    When the formulas are representable in a folded scheme, we find tensor network decompositions of $\fselectionmap$ and exploit them along efficient representations of $\kb$ in an efficient calculation of $\contraction{\kb,\sencodingof{\formulaset}} $.
    This is further equal to
    \[ \kb \models \lnot \left( \bigvee_{\exformula\in\formulaset} \exformula\right) \, . \]
\end{example}
